\chapter*{How to read this document}
\addcontentsline{toc}{chapter}{How to read this document}

This thesis is written to be read in two ways, and it is separated so that either
works.

\paragraph{The argument} is Parts~I--IV, Chapters~\ref{ch:intro}--\ref{ch:plan}.
It assumes you know what a photograph is and are willing to look at a
$3\times3$ matrix, and nothing beyond that. It contains the question, the
derivation, what was measured, and what the measurements mean. Every chapter
opens with the idea in plain language before any mathematics, and no chapter in
this part contains a file path, a command line, or a discussion of how the
computation was arranged. Read straight through, it is a talk.

\paragraph{The technical record} is the appendices. Appendix~\ref{app:derivations}
carries the derivations in full, for a reader who would sooner check the algebra
than accept the results. Appendix~\ref{app:harness} describes the machinery
that produced the measurements --- where the work ran, what is asserted before
anything is measured, how a result enters the record and how it leaves it.
Appendix~\ref{app:impl} gives the implementation: four configurations of one codebase, and what each switch touches. Appendix~\ref{app:defects} is the
eight defects found in this project's own measurement, each with what it looked
like, what was actually happening, and what now prevents it.

\medskip
\noindent
That last appendix is the one a supervisor is most likely to want and a seminar
audience is least likely to. It is kept complete and kept out of the way.

\section*{If you have twenty minutes}

Read the Abstract, then \S\ref{sec:table1} and its figures --- the reproduction
and what a ten-decibel gap looks like as a picture --- then \S\ref{sec:table3},
the experiment the thesis exists to perform, then Chapter~\ref{ch:discussion} for
what it does and does not establish. Four tables and three figures carry the
whole argument.

\section*{A note on what is claimed}

This document reports a negative result and a null result alongside its positive
one, and it does not weight them differently. Method~II is refuted. Method~I is
exactly equal to its baseline on every standard benchmark --- by construction,
provably, and that is the intended behaviour, not a disappointment. Its
one positive result is $0.197$\,dB on a single degraded capture protocol, at a
cost of $2.8\times$ the memory.

The reason to read on is not that the method works. It is that the derivation
predicted, in advance and in detail, \emph{where} it would do nothing and where
it would act, and the measurement agrees --- including on the five protocols
where the prediction was that nothing would happen. A theory that correctly
predicts its own null results is doing more work than one that only predicts its
successes.
